\section{开元的名册、粮道与宫廷边界} \label{sec:xuanzong-households-transport-court}  710年至713年的政变、受禅和清洗结束了并存的权力中心，却没有让国家机器从此自行运转。玄宗掌握中枢后，仍须决定哪些户进入赋役名册、哪些粮食沿何种水路输送、哪些官员能够在京师和地方之间传递命令。把713年视为开元秩序的自然起点，会遮住这种靠禁军、宗室和文官网络取得的政治前提。% \footnote{唐隆政变、玄宗即位与太平公主集团覆灭见 ST-710-TANGLONG-COUP、ST-712-XUANZONG-ACCESSION、ST-713-TAIPING，依据《旧唐书》《新唐书》睿宗、玄宗纪及《资治通鉴》。胜利者叙事不足以完整复原政变中的命令、犹疑与未参与者。}  721年宇文融括户检田，把逃户和隐漏户重新纳入登记；它既是增加赋役资源的政策，也是国家承认原有名册未能覆盖实际人口与土地的一刻。登记在文书中显得整齐，落到村落却意味着身份、居住、田地与负担的重新认定。我们不能从括得的数字推算全国人口，更不能把被登记者都写成被动对象；但地方文书的零星保存足以提醒我们，国家看到家庭的方式与家庭理解自身生活的方式并不相同。% \footnote{宇文融括户检田见 ST-721-YUWEN-RONG，依据《旧唐书·宇文融传》《新唐书·食货志》《通典》及可与之对照的敦煌、吐鲁番户籍文书。不同地区和年份的登记口径不可直接合并。}  733年裴耀卿整顿江淮漕运，设置转运节点以改善关中供粮。漕运的意义不只在于让长安“富足”：仓、船、河道、役夫、商人和各级官员必须在时间上接合，任一环失灵都会将压力转回地方。运河与仓储遗存可以帮助辨认路线和维护，却不能据此精算运量或征役人数；正史中的成效叙述也不该取代沿线居民承担的土地、劳力和市场波动。% \footnote{裴耀卿漕运见 ST-733-PEI-YAOQING，依据《旧唐书·裴耀卿传》《新唐书·食货志》《通典》及运河、仓储考古资料。制度性总结和工程遗存都不能直接推出单年运输总量。}  725年封禅、736年李林甫拜相、737年三王被废而李亨继立，展示出开元晚期的礼仪成功、相权集中和继承不安同时发生。742年改州为郡、刺史为太守，是一项重新整理行政语言的尝试；名称的整齐却不能保证财政、兵役和地方命令也同样整齐。天宝以前的繁荣应当被理解为一种需要不断维修的连接，而不是能从封禅、户口或官名中一次读出的社会共识。% \footnote{封禅、李林甫、三王废立及天宝改制见 ST-725-TAISHAN、ST-736-LI-LINFU、ST-737-CROWN-PRINCES、ST-742-TIANBAO，依据两《唐书》玄宗纪、李林甫传、礼仪和地理志及《资治通鉴》。人物责任和地方执行均不能以安史结局倒推。}